
\documentclass{article}
\usepackage{colortbl}
\usepackage{makecell}
\usepackage{multirow}
\usepackage{supertabular}

\begin{document}

\newcounter{utterance}

\centering \large Interaction Transcript for game `wordle\_withcritic', experiment `medium\_frequency\_words\_clue\_with\_critic', episode 2 with reka{-}flash{-}3.1{-}t0.0.
\vspace{24pt}

{ \footnotesize  \setcounter{utterance}{1}
\setlength{\tabcolsep}{0pt}
\begin{supertabular}{c@{$\;$}|p{.15\linewidth}@{}p{.15\linewidth}p{.15\linewidth}p{.15\linewidth}p{.15\linewidth}p{.15\linewidth}}
    \# & $\;$A & \multicolumn{4}{c}{Game Master} & $\;\:$B\\
    \hline

    \theutterance \stepcounter{utterance}  
    & & \multicolumn{4}{p{0.6\linewidth}}{
        \cellcolor[rgb]{0.9,0.9,0.9}{
            \makecell[{{p{\linewidth}}}]{
                \texttt{\tiny{[P1$\langle$GM]}}
                \texttt{You are a language wizard who likes to guess words by using the given rules.} \\
\\ 
\texttt{Welcome to Wordle! You have six attempts to guess the target word, a valid English word of five lowercase letters (a{-}z). Please use the tags "explanation:" and "guess:" to provide a concise explanation for each guess.} \\
\\ 
\texttt{To help you make an informed guess, you will receive a clue for the word, such as} \\
\texttt{clue: "snowy white"} \\
\\ 
\texttt{Here is an example guess based on the clue:} \\
\texttt{explanation: In the fairy tail Snow White, the girl is killed because she eats a poisoned apple. And the word apple has 5 letters.} \\
\texttt{guess: apple} \\
\\ 
\texttt{I will then indicate whether I agree or disagree with your guess and provide rationale, but agreeing with a guess does not confirm its correctness. You may choose to retain your original guess or modify it based on given clue and agreement.} \\
\\ 
\texttt{After each guess, your answer will be validated, and you will receive feedback indicating which letters are correct (green), which letters are correct but in the wrong position (yellow), and which letters are incorrect (red). This feedback can be useful in determining which letters to include or exclude in your next guess.} \\
\\ 
\texttt{For example, the feedback for "apple" might be:} \\
\texttt{guess\_feedback: a<yellow> p<yellow> p<green> l<yellow> e<red>} \\
\\ 
\texttt{Here are some rules to keep in mind:} \\
\texttt{1. Please do not generate any guess feedback or clue on your own} \\
\texttt{2. Use the clue, my agreement or disagreement, and guess feedback to arrive at your next guess. This should be reflected in your explanation as well.} \\
\\ 
\texttt{Do not generate any other text.} \\
\\ 
\texttt{Let's begin with your first guess.} \\
\\ 
\texttt{clue: plastics base} \\
            }
        }
    }
    & \\ \\

    \theutterance \stepcounter{utterance}  
    & \multicolumn{4}{p{0.6\linewidth}}{
        \cellcolor[rgb]{0.9,0.9,0.9}{
            \makecell[{{p{\linewidth}}}]{
                \texttt{\tiny{[P1$\rangle$GM]}}
                \texttt{<reasoning>} \\
\texttt{Let's tackle this Wordle with the clue "plastics base". First, I need to understand what the clue is pointing towards. The words "plastics" and "base" are the key here. Maybe the target word is related to both plastics and a base? Or perhaps it's a compound word or a term that combines elements from both.} \\
\\ 
\texttt{Plastics are synthetic materials, often used in manufacturing. A base could refer to a foundation, a sports base, or maybe a chemical base. Let me think of words that connect these two. Maybe "plastic" has a root that's part of the word. "Plastic" comes from the Greek word "plastos" meaning molded. But how does that connect to base?} \\
\\ 
\texttt{Alternatively, maybe the clue is a play on words. "Plastics base" could be a pun or a homophone. Let me think of words that sound like parts of these. "Plastics" could be split into "plat" and "ics", but not sure. "Base" could be "bass" or "baze". Not obvious.} \\
\\ 
\texttt{Wait, maybe the target word is "plastic" itself, but that's six letters. The word needs to be five letters. So maybe a shorter form? Or perhaps a word that incorporates "base" and relates to plastics. Let's think of five{-}letter words. Words like "bases" (plural of base) but that's five letters. Does "bases" relate to plastics? Not directly. Maybe "plast" is part of it, but "plast" is four letters. } \\
\\ 
\texttt{Alternatively, think of common five{-}letter words that are associated with both plastics and a base. Maybe "molded" but that's six letters. "Mold" is five letters. Does "mold" relate to plastics? Plastics can be molded, so maybe. But how does "base" come into play? Maybe "mold" is part of a base structure? Not sure.} \\
\\ 
\texttt{Another angle: "plastics base" could be a term in a specific context. For example, in construction, a base layer made of plastics? Or in chemistry, a base used in plastic manufacturing? Maybe "alkali" is a base used in plastics, but that's six letters. "Sodium" is too long. } \\
\\ 
\texttt{Wait, maybe the clue is a riddle. "Plastics base" – maybe the word is "plaste} \\
            }
        }
    }
    & & \\ \\

    \theutterance \stepcounter{utterance}  
    & & & \multicolumn{2}{p{0.3\linewidth}}{
        \cellcolor[rgb]{0.9,0.9,0.9}{
            \makecell[{{p{\linewidth}}}]{
                \texttt{\tiny{[GM$|$GM]}}
                \texttt{Error: The response should always start with the keyword 'explanation:'} \\
            }
        }
    }
    & & \\ \\

    \theutterance \stepcounter{utterance}  
    & & & \multicolumn{2}{p{0.3\linewidth}}{
        \cellcolor[rgb]{0.9,0.9,0.9}{
            \makecell[{{p{\linewidth}}}]{
                \texttt{\tiny{[GM$|$GM]}}
                \texttt{game\_result = ABORT} \\
            }
        }
    }
    & & \\ \\

\end{supertabular}
}

\end{document}
